\chapter{正负数的乘法}\label{ch:g8:negprod}

七年级学习了\omterm{def:g7:negatives:def}{正负数}的加减法
(\cref{ch:g7:negatives})，还需学习它们的乘法和除法。本章
的核心问题是：为什么“负乘负”应当得正？
我们给出理由，而不只是给出法则。

\section{符号法则}

\begin{theorem}[积的符号]\label{thm:g8:negprod:rules}
两个\omterm{def:g7:negatives:def}{正负数}的积的绝对值
等于绝对值之积，其符号由下表确定：
\begin{center}
\renewcommand{\arraystretch}{1.3}
\begin{tabular}{c|cc}
$\times$ & $+$ & $-$ \\
\hline
$+$ & $+$ & $-$ \\
$-$ & $-$ & $+$ \\
\end{tabular}
\end{center}
\emph{同号得正，异号得负。}
商的符号同样由此表确定。
\end{theorem}

\begin{proof}[为什么 $(-1) \times (-1) = +1$]
首先，$3 \times (-4)$ 表示 $(-4) + (-4) + (-4) = -12$：正数
乘负数的结果为负。现观察下面的规律：
\[
3 \times (-4) = -12, \quad
2 \times (-4) = -8, \quad
1 \times (-4) = -4, \quad
0 \times (-4) = 0 .
\]
每一步中第一个\omterm{def:g5:division:multiple}{因数}减少 $1$，结果\emph{增大}
增加 $4$。沿此规律再进行一步：
\[
(-1) \times (-4) = 0 + 4 = +4 .
\]
若选取其他结果，就会破坏算术运算的规律性
（分配律）。因此，负乘负必得正。
\end{proof}

\begin{omfigure}
\begin{tikzpicture}
\begin{axis}[omaxis,
  width=7.8cm, height=5.4cm,
  xmin=-2.6, xmax=3.6, ymin=-14, ymax=10,
  xlabel={$k$}, ylabel={$k \times (-4)$},
  xtick={-2,-1,1,2,3}, ytick={-12,-8,-4,4,8}]
  \addplot[omDef, thick, domain=-2.3:3.3] {-4*x};
  \addplot[omDef, only marks, mark=*, mark size=1.6pt] coordinates
    {(3,-12) (2,-8) (1,-4) (0,0) (-1,4) (-2,8)};
\end{axis}
\end{tikzpicture}

{\small 用图形理解规律：各点
$(k,\ k \times (-4))$ 排成一条直线，将该直线延伸越过 $k = 0$，
可得 $(-1)\times(-4) = +4$ 且 $(-2)\times(-4) = +8$。}
\end{omfigure}

\begin{example}\label{ex:g8:negprod:products}
\begin{align*}
(-7) \times (+6) &= -42 && \text{（异号）} \\
(-5) \times (-8) &= +40 && \text{（同号）} \\
(+2.5) \times (-4) &= -10 \\
(-36) \div (-9) &= +4 \\
(+15) \div (-2) &= -7.5 .
\end{align*}
\end{example}

\section{多个因数的积}

\begin{proposition}[多个因数之积的符号]\label{prop:g8:negprod:many}
若干个非零\omterm{def:g5:division:multiple}{因数}之积的符号只取决于
负\omterm{def:g5:division:multiple}{因数}的个数：
\begin{itemize}
  \item 负\omterm{def:g5:division:multiple}{因数}的个数为\emph{\omterm{def:g3:numbers:evenodd}{偶数}}时，
        积为正；
  \item 负\omterm{def:g5:division:multiple}{因数}的个数为\emph{\omterm{def:g3:numbers:evenodd}{奇数}}时，积为负。
\end{itemize}
\end{proposition}

\begin{proof}
负\omterm{def:g5:division:multiple}{因数}两两配对，每一对都给出正号
（由\cref{thm:g8:negprod:rules}可知）；若负\omterm{def:g5:division:multiple}{因数}个数为\omterm{def:g3:numbers:evenodd}{奇数}，则
剩下一个未配对的负\omterm{def:g5:division:multiple}{因数}，使整个积为负。
\end{proof}

\begin{example}\label{ex:g8:negprod:many}
$(-2) \times (+3) \times (-5) \times (-1)$：三个负\omterm{def:g5:division:multiple}{因数}
（\omterm{def:g3:numbers:evenodd}{奇数}个），故积为负；绝对值之积为 $2 \times 3 \times 5
\times 1 = 30$。结果为 $-30$。先\emph{确定符号}，再
计算绝对值——
两个简单的任务代替一个易错的任务。
\end{example}

\begin{example}[负数的幂]\label{ex:g8:negprod:powers}
$(-2)^2 = (-2) \times (-2) = +4$，且
$(-2)^3 = (-2)^2 \times (-2) = -8$：\omterm{def:g3:numbers:evenodd}{偶数}次幂的
值为正，\omterm{def:g3:numbers:evenodd}{奇数}次幂保持底数的符号。注意书写：
$(-3)^2 = 9$，而 $-3^2 = -(3 \times 3) = -9$。
\end{example}

\section{四则混合运算}

\begin{method}[正负数计算的稳妥方法]\label{met:g8:negprod:compute}
\begin{enumerate}
  \item 遵循\cref{ch:g7:priorities}中的运算顺序：先算括号，
        再算 $\times$ 和 $\div$，最后算 $+$ 和 $-$；
  \item 每次做乘法或除法时，先求\emph{符号}，
        再求绝对值；
  \item 每行只写一步。
\end{enumerate}
\end{method}

\begin{example}\label{ex:g8:negprod:compute}
\begin{align*}
5 - 3 \times (-4)
&= 5 - (-12) && \text{（先算乘法：异号）}\\
&= 5 + 12 = 17 .
\end{align*}
\begin{align*}
\frac{(-6) \times (-10)}{-4}
&= \frac{+60}{-4} && \text{（分子：同号）}\\
&= -15 && \text{（商：异号）。}
\end{align*}
\end{example}

\begin{example}[代入负数求值]\label{ex:g8:negprod:substitute}
求 $E = 3x^2 - 2x$ 当 $x = -5$ 时的值，代入时数值须加括号：
\[
E = 3 \times (-5)^2 - 2 \times (-5)
= 3 \times 25 - (-10)
= 75 + 10 = 85 .
\]
$(-5)^2$ 中的括号必不可少：若无括号，
则平方只会作用于 $5$。
\end{example}

\section{练习}

\begin{exercise}[$\star$]\label{exo:g8:negprod:1}
计算：
\[
(-8) \times (+7), \qquad
(-9) \times (-6), \qquad
(+12) \times (-0.5), \qquad
(-1) \times (-1) \times (-1).
\]
\end{exercise}

\begin{exercise}[$\star$]\label{exo:g8:negprod:2}
计算：
\[
(-63) \div (+9), \qquad
(-8) \div (-16), \qquad
\frac{45}{-5}, \qquad
\frac{-7.2}{-8} .
\]
\end{exercise}

\begin{exercise}[$\star$]\label{exo:g8:negprod:3}
只写出每个乘积的\emph{符号}，无需计算：
\[
(-3) \times (-7) \times (+2) \times (-5); \qquad
(-1)^{10}; \qquad
(-2)^{7}; \qquad
(-4) \times 0 \times (-6).
\]
\end{exercise}

\begin{exercise}[$\star$]\label{exo:g8:negprod:4}
计算：
\[
(-2)^4, \qquad
-2^4, \qquad
(-3)^3, \qquad
(-1)^{2026} .
\]
\end{exercise}

\begin{exercise}[$\star$]\label{exo:g8:negprod:5}
遵循运算顺序，逐步计算：
\[
7 + 2 \times (-6), \qquad
(-4) \times (5 - 9), \qquad
18 \div (-3) - (-2) .
\]
\end{exercise}

\begin{exercise}[$\star$]\label{exo:g8:negprod:6}
计算：
\[
\frac{(-5) \times (+8)}{-4}, \qquad
\frac{(-3) \times (-6)}{(-2) \times (+9)} .
\]
\end{exercise}

\begin{exercise}[$\star$]\label{exo:g8:negprod:7}
求 $x = -3$ 时下列各式的值：
\[
5x, \qquad x^2, \qquad -x^2, \qquad 2x^2 + 4x, \qquad (2x)^2 .
\]
\end{exercise}

\begin{exercise}[$\star\star$]\label{exo:g8:negprod:8}
补全下列各等式：
\[
(-6) \times \;?\; = 42, \qquad
\;?\; \div (-4) = -2.5, \qquad
(-5) \times \;?\; = -1 .
\]
\end{exercise}

\begin{exercise}[$\star\star$]\label{exo:g8:negprod:9}
判断下列说法是否正确，并说明理由或举出反例。
\begin{enumerate}
  \item 两个\omterm{def:g7:negatives:def}{正负数}的积总不小于它们的
        和。
  \item \omterm{def:g7:negatives:def}{正负数}的平方永不为负。
  \item 若三个\omterm{def:g5:division:multiple}{因数}的积为正，则这三个\omterm{def:g5:division:multiple}{因数}
        全为正数。
\end{enumerate}
\end{exercise}

\begin{exercise}[$\star\star$]\label{exo:g8:negprod:10}
测验中每道答错的题记 $-3$ 分，每道答对的
题记 $+5$ 分。佐伊答了 $20$ 道题，得 $36$ 分。
答对的题有几道？（可试取值，或列出计算式
$5r - 3(20 - r) = 36$。）
\end{exercise}

\begin{exercise}[$\star\star\star$]\label{exo:g8:negprod:11}
利用分配律（与证明
\cref{thm:g8:negprod:rules}时一样），展开并论证下式的每一步
\[
0 = (-1) \times 0 = (-1) \times \bigl(1 + (-1)\bigr)
= (-1) \times 1 + (-1) \times (-1),
\]
并由此得出 $(-1) \times (-1)$ 必等于 $+1$。
\end{exercise}

\section{问题：为什么负负得正}

\begin{problem}[{周末问题——由
分配律推出的符号法则，以及交错和 $1 - 2 + 3 - 4 + \dots$}]
\label{pb:g8:negprod:1}
\cref{thm:g8:negprod:rules}的证明延续了一种规律，并断言
“任何其他选择都会破坏分配律”。本题
将使这一论断精确化：仅从分配律出发
（见\cref{thm:g7:priorities:distributivity}），你将\emph{证明}
符号法则——不靠图形，不靠找规律，而是按代数本身的方式——
然后将其用于 $-1$ 的幂，以及一个著名的
交错和式。全文中 $a$ 与 $b$ 为\omterm{def:g7:negatives:def}{正负数}，
并默认\cref{ch:g7:negatives}中的加法法则，以及
$1 \times a = a$，并默认乘积可按任意次序计算
（因此分配律对乘积的两侧均适用）。

\medskip
\noindent\textbf{第一部分——符号法则是定理。}
\begin{enumerate}
  \item 证明 $0 \times a = 0$ 对每个\omterm{def:g7:negatives:def}{正负数} $a$ 成立。
        （写出 $0 = 0 + 0$，展开 $(0 + 0) \times a$，并问：
        哪个数加上 $0 \times a$ 后仍得 $0 \times a$？）
  \item 展开 $\bigl(1 + (-1)\bigr) \times a$，以证明
        $(-1) \times a + a = 0$，并由此推出
        \[
        (-1) \times a = -a :
        \]
        乘以 $-1$ 即得\omterm{def:g7:negatives:def}{相反数}。可与
        \cref{exo:g8:negprod:11}比较，它正是 $a = -1$ 的情形。
  \item 推出 $(-a) \times b = -(a \times b)$：一个负号
        可原样提到乘积之外。
  \item 推出 $(-a) \times (-b) = a \times b$：两个负号
        相互抵消。
  \item 每个\omterm{def:g7:negatives:def}{正负数}都是其到零的距离前冠以一个符号：
        $a = +d$ 或 $a = -d$。说明第 3、4 题如何
        证明\cref{thm:g8:negprod:rules}中符号表的全部
        四种情形，连同距离一并给出。
\end{enumerate}

\medskip
\noindent\textbf{第二部分——$-1$ 的幂。}
\begin{enumerate}[resume]
  \item 计算 $(-1)^2$、$(-1)^3$、$(-1)^4$、$(-1)^5$，然后给出
        依据\cref{prop:g8:negprod:many}，给出
        $(-1)^n$ 的值，其中 $n \geq 1$ 为正整数。
  \item 不计算任何距离，给出下式的符号：
        \[
        (-1) \times (-2) \times (-3) \times \dots \times (-10),
        \]
        并把此积的距离写成正整数的乘积，
        无需算出其值。
  \item 证明 $(-a)^2 = a^2$，且 $(-a)^3 = -a^3$。负号在哪些
        指数下得以保留？
  \item 利用符号法则证明，\omterm{def:g7:negatives:def}{正负数}的平方
        永不为负，并由此推出没有\omterm{def:g7:negatives:def}{正负数}
        $x$ 满足 $x^2 = -4$。
  \item 佐伊断言：“$(-1)^{2026} + (-1)^{2027} = 0$，更一般地，
        $-1$ 的两个相邻幂总相互抵消。”她说得对吗？
        请说明理由。
\end{enumerate}

\medskip
\noindent\textbf{第三部分——交错和。}
符号法则既已确立，便可计算兼含两种符号的和。
对正整数 $n \geq 1$，令
\[
A_n = 1 - 2 + 3 - 4 + \dots
\]
为从 $1$ 到 $n$ 的各正整数依交错符号相加之和，
末项为 $+n$（当 $n$ 为\omterm{def:g3:numbers:evenodd}{奇数}时），$-n$（当 $n$ 为\omterm{def:g3:numbers:evenodd}{偶数}时）。
于是 $A_1 = 1$，$A_2 = 1 - 2 = -1$。
\begin{enumerate}[resume]
  \item 计算 $A_3$、$A_4$、$A_5$、$A_6$ 与 $A_7$。
        你对此有何猜想？
  \item 设 $n$ 为\omterm{def:g3:numbers:evenodd}{偶数}。将各项两两分组，
        得 $(1 - 2) + (3 - 4) + \dots$，数出其中的组数，并证明
        $A_n = -\frac{n}{2}$。
  \item 设 $n$ 为\omterm{def:g3:numbers:evenodd}{奇数}。说明为何 $A_n = A_{n-1} + n$，并
        由上一小题推出
        $A_n = \frac{n + 1}{2}$。
  \item 计算 $A_{100}$、$A_{101}$ 与 $A_{2026}$。
  \item 阿尔玛算得 $A_{101} - A_{100} = 51 - (-50) = 101$，并
        惊讶地发现结果恰为 $101$。她本该惊讶吗？
        请用一句话解释，并不经任何进一步计算写出
        $A_{2027} - A_{2026}$ 的值。
\end{enumerate}
\end{problem}
